% ============================================================================
% Slide De tai 09: MCU - SoC - FPGA. Bien dich: xelatex slides.tex (2 lan)
% ============================================================================
\documentclass[aspectratio=169,11pt]{beamer}
% ============================================================================
% Phan dung chung cho report.tex va slides.tex (bien dich bang XeLaTeX)
% ============================================================================
\usepackage{fontspec}
\usepackage{polyglossia}
\setdefaultlanguage{vietnamese}
\setotherlanguage{english}
\setmonofont{DejaVu Sans Mono}[Scale=0.85]
\usepackage{amsmath}
\usepackage{newunicodechar}
\newunicodechar{≈}{\ensuremath{\approx}}
\newunicodechar{−}{\ensuremath{-}}

\usepackage{xcolor}
\usepackage{graphicx}
\usepackage{booktabs}
\usepackage{tabularx}
\usepackage{array}
\usepackage{makecell}

% Mau theo nen tang (giong chart trong experiments/results/charts)
\definecolor{cMCU}{HTML}{2A78D6}
\definecolor{cSoC}{HTML}{EB6834}
\definecolor{cFPGA}{HTML}{1BAF7A}
\definecolor{cInk2}{HTML}{52514E}
\definecolor{cWarn}{HTML}{B54708}

% Nhan nguon goc so lieu: \nhan{do}, \nhan{mo phong}, ...
\newcommand{\nhan}[1]{\textcolor{cInk2}{\footnotesize[#1]}}
% So chua do: hien ro rang, khong bao gio dien so gia dinh
\newcommand{\chuado}{\textcolor{cWarn}{\textit{chưa đo}}}

\newcommand{\MCU}{\textcolor{cMCU}{\textbf{MCU}}}
\newcommand{\SoC}{\textcolor{cSoC}{\textbf{SoC}}}
\newcommand{\FPGA}{\textcolor{cFPGA}{\textbf{FPGA}}}

% So lieu thi nghiem, tu dong sinh boi experiments/e5_crosscheck/collect.py
\newcommand{\resdir}{../experiments/results}
\IfFileExists{\resdir/numbers.tex}{\input{\resdir/numbers.tex}}{%
  \PackageWarning{shared}{Chua co numbers.tex - chay: make -C ../experiments collect}%
  % Ban du phong khi chua chay collect.py: moi so deu \chuado{}
% So chua do hien \chuado{} (report dinh nghia macro nay).
\newcommand{\expMcuClock}{\chuado{}}
\newcommand{\expMcuBestVariant}{\chuado{}}
\newcommand{\expMcuBestCps}{\chuado{}}
\newcommand{\expMcuBestMsps}{\chuado{}}
\newcommand{\expMcuEff}{\chuado{}}
\newcommand{\expMcuShiftCps}{\chuado{}}
\newcommand{\expMcuIrqMin}{\chuado{}}
\newcommand{\expMcuIrqMax}{\chuado{}}
\newcommand{\expMcuIrqRange}{\chuado{}}
\newcommand{\expMcuIrqHogMax}{\chuado{}}
\newcommand{\expPiBestMsps}{\chuado{}}
\newcommand{\expFpgaMsps}{\chuado{}}
\newcommand{\expFpgaLabel}{\chuado{}}
\newcommand{\expFpgaLatCycles}{\chuado{}}
\newcommand{\expFpgaLatNs}{\chuado{}}
\newcommand{\expFpgaPass}{\chuado{}}
\newcommand{\expFpgaSpc}{\chuado{}}
\newcommand{\expFpgaFmax}{\chuado{}}
\newcommand{\expFpgaDsp}{\chuado{}}
\newcommand{\expFpgaPower}{\chuado{}}
\newcommand{\expRatioFpgaPi}{\chuado{}}
\newcommand{\expRatioFpgaMcu}{\chuado{}}
\newcommand{\expLinuxStdMax}{\chuado{}}
\newcommand{\expLinuxRtMax}{\chuado{}}
\newcommand{\expHaveMcu}{0}
\newcommand{\expHavePi}{0}
}

% Chen chart neu da co, neu chua thi hien khung bao "chua co"
\newcommand{\chart}[2][\linewidth]{%
  \IfFileExists{\resdir/charts/#2}%
    {\includegraphics[width=#1]{\resdir/charts/#2}}%
    {\fbox{\parbox{0.8\linewidth}{\centering\chuado{}: chưa có \texttt{#2}}}}}

\setsansfont{Noto Sans}
\setmainfont{Noto Sans}

\usetheme{default}
\usefonttheme{professionalfonts}
\setbeamertemplate{navigation symbols}{}
\setbeamertemplate{footline}{%
  \hfill\textcolor{cInk2}{\scriptsize CO3053 · Đề tài 09 \quad \insertframenumber/\inserttotalframenumber}\hspace*{1em}\vspace*{0.6em}}
\setbeamercolor{frametitle}{fg=black}
\setbeamerfont{frametitle}{series=\bfseries,size=\Large}
\setbeamercolor{structure}{fg=cInk2}
\setbeamertemplate{itemize items}{\textbullet}
\setbeamersize{text margin left=1.2cm,text margin right=1.2cm}
\renewcommand{\arraystretch}{1.25}
\newcolumntype{Y}{>{\raggedright\arraybackslash}X}

% Khoi "con so quyet dinh"
\newcommand{\soquyetdinh}[1]{%
  \begin{beamercolorbox}[rounded=true,sep=4pt]{block body}\footnotesize #1\end{beamercolorbox}\vspace{4pt}}
\setbeamercolor{block body}{bg=black!4}

\title{\textbf{MCU – SoC – FPGA}}
\subtitle{Đánh đổi giữa hiệu năng, chi phí và độ linh hoạt}
\author{CO3053 Hệ thống nhúng · HK261 · Đề tài 09}
\date{}

\begin{document}

% ---------------------------------------------------------------- 1
\begin{frame}[plain]
\vfill
{\Huge\bfseries \MCU{} – \SoC{} – \FPGA}\par\medskip
{\Large Đánh đổi giữa hiệu năng, chi phí và độ linh hoạt}\par\bigskip
{\color{cInk2} CO3053 Hệ thống nhúng · HK261 · Đề tài 09}
\vfill
\end{frame}

% ---------------------------------------------------------------- 2
\begin{frame}{Câu hỏi của bài}
\begin{center}
{\Large Khi nào nên chọn \MCU, \SoC{} hay \FPGA?}
\end{center}
\bigskip
\begin{itemize}
\item Xét theo 3 tiêu chí: \textbf{hiệu năng}, \textbf{chi phí}, \textbf{độ linh hoạt}
\item \emph{Không} nhằm chứng minh nền tảng nào ``tốt nhất'', vì đó là mục tiêu thiên vị mà \texttt{pe2} cảnh báo
\item Mỗi con số trên slide có một nhãn cho biết nguồn gốc:\\
      \nhan{đo} \nhan{mô phỏng} \nhan{ước lượng công cụ} \nhan{tính} \nhan{nguồn}
\end{itemize}
\end{frame}

% ---------------------------------------------------------------- 3
\begin{frame}{Ba kiến trúc}
\begin{columns}[T]
\begin{column}{0.32\linewidth}
\MCU\par\smallskip\small
Cortex-M, flash và SRAM trên chip\\[4pt]
Chạy lệnh tuần tự\\
\textbf{Không có MMU}, nên không chạy Linux\\
Khởi động trong vài ms, ngắt dưới 1 µs
\end{column}
\begin{column}{0.32\linewidth}
\SoC\par\smallskip\small
Cortex-A có MMU và cache, cùng GPU, codec, khối tăng tốc mạng\\[4pt]
Chạy \textbf{Linux}: driver, mạng, file system có sẵn\\
Đổi lại: độ trễ khó đoán trước
\end{column}
\begin{column}{0.32\linewidth}
\FPGA\par\smallskip\small
Mảng LUT, DSP, BRAM nối tùy ý\\[4pt]
Thuật toán ``trải ra không gian'': mỗi phép nhân có một khối riêng\\
1 mẫu/chu kỳ, độ trễ cố định\\
Đổi lại: phải tự mô tả mọi thứ bằng HDL
\end{column}
\end{columns}
\vfill
\centering\small\textcolor{cInk2}{Zynq (CPU và FPGA trên một chip) được xếp vào \textbf{phương án lai}}
\end{frame}

% ---------------------------------------------------------------- 4
\begin{frame}{Số liệu từ tài liệu: hiệu năng}
\begin{columns}[T]
\begin{column}{0.48\linewidth}
\small\textbf{CoreMark} \nhan{nguồn}\par\smallskip
\begin{tabular}{@{}lr@{}}
\toprule
\MCU{} STM32F4 (M4, 180 MHz) & 608\\
\MCU{} STM32H743 (M7, 480 MHz) & 2\,424\\
\SoC{} STM32MP1 (A7, 800 MHz) & 4\,158\\
\SoC{} Pi 4 (4×A72, 1,8 GHz) & 38\,200\\
\bottomrule
\end{tabular}
\par\smallskip\scriptsize\textcolor{cInk2}{CoreMark không áp dụng cho FPGA}
\end{column}
\begin{column}{0.48\linewidth}
\small\textbf{Lọc video 1080p 30 khung/giây cần ≈1,68 G MAC/s}\par\smallskip
\begin{tabular}{@{}lr@{}}
\toprule
\MCU{} H743 & ≈0,48 G \nhan{tính}\\
\FPGA{} XC7A100T, 240 DSP & ≈48 G \nhan{tính}\\
\bottomrule
\end{tabular}
\par\medskip
FIR trên Zynq nhanh hơn Cortex-A9 \textbf{≈27×} \nhan{nguồn}\par
và tốn ít năng lượng hơn ≈95–190× \nhan{tính}
\end{column}
\end{columns}
\end{frame}

% ---------------------------------------------------------------- 5
\begin{frame}{Số liệu từ tài liệu: độ trễ, giá, công suất}
\begin{columns}[T]
\begin{column}{0.46\linewidth}
\small\textbf{Độ trễ xấu nhất} \nhan{nguồn: AC6, STAC}\par\smallskip
\begin{tabular}{@{}lr@{}}
\toprule
\FPGA{} tick-to-trade & 13,9 ns\\
\MCU{} bare-metal & 1–3 µs\\
\SoC{} PREEMPT\_RT & 30–80 µs\\
\SoC{} Linux thường & vài ms\\
\bottomrule
\end{tabular}
\end{column}
\begin{column}{0.50\linewidth}
\small\textbf{Giá và công suất} \nhan{nguồn}\par\smallskip
\begin{tabular}{@{}lll@{}}
\toprule
& Giá & Khi nghỉ\\
\midrule
\MCU & £9–12/chip & 300 nA (U5)\\
\SoC & \$35/board & ≈3 W\\
\FPGA & €144/chip & chưa có số\\
\bottomrule
\end{tabular}
\par\smallskip\scriptsize\textcolor{cInk2}{So chip với board: chỉ dùng để so bậc độ lớn}
\end{column}
\end{columns}
\end{frame}

% ---------------------------------------------------------------- 6
\begin{frame}{Thí nghiệm: một bài toán, ba nền tảng}
\textbf{FIR 16-tap Q15}, 1024 mẫu: bậc thang, giá trị biên, mẫu ép bão hòa, mẫu ngẫu nhiên\par
Mọi nền tảng phải khớp \textbf{từng mẫu} với golden model viết bằng Python\par\medskip
\small
\begin{tabularx}{\linewidth}{@{}l l Y@{}}
\toprule
\MCU  & STM32F103, Cortex-M3 & 4 phiên bản (C, CMSIS-DSP), đếm chu kỳ bằng DWT; độ trễ ngắt trong 3 kịch bản\\
\SoC  & Raspberry Pi 4 & Cùng mã C, \texttt{-O2} và \texttt{-O3} (NEON); \texttt{cyclictest} khi rảnh và khi có tải\\
\FPGA & Arty A7-100T & Pipeline 7 tầng; tự so với golden ngay trên chip ở 100 MHz; Vivado cho Fmax và công suất\\
\bottomrule
\end{tabularx}
\par\medskip
\textcolor{cInk2}{Cortex-M3 không có lệnh SIMD (SMLAD), nên CMSIS-DSP không được tăng tốc kiểu SIMD như trên M4}
\end{frame}

% ---------------------------------------------------------------- 7
\begin{frame}{Kết quả: thông lượng}
\centering\chart[0.92\linewidth]{e5_throughput.png}
\par\small
\FPGA: \expFpgaSpc{} mẫu/chu kỳ, \expFpgaMsps{} MSPS \nhan{\expFpgaLabel} \quad
\MCU: \expMcuBestMsps{} MSPS \nhan{đo} \quad
\SoC: \expPiBestMsps{} MSPS \nhan{đo}
\end{frame}

% ---------------------------------------------------------------- 8
\begin{frame}{Kết quả: độ trễ xấu nhất}
\centering\chart[0.92\linewidth]{e5_worst_latency.png}
\par\small\raggedright
\MCU{} ngắt khi chạy FIR: \expMcuIrqRange{} ns · thêm ngắt ưu tiên cao hơn 20 µs: \expMcuIrqHogMax{} µs \nhan{đo}\\
\SoC{} Linux thường: \expLinuxStdMax{} µs · PREEMPT\_RT: \expLinuxRtMax{} µs \nhan{đo} \quad
\FPGA{} pipeline: \expFpgaLatCycles{} chu kỳ = \expFpgaLatNs{} ns \nhan{\expFpgaLabel}
\end{frame}

% ---------------------------------------------------------------- 9
\begin{frame}{Kiểm chứng chéo với tài liệu}
\small
\begin{tabularx}{\linewidth}{@{}Y Y Y@{}}
\toprule
Câu hỏi & Nhóm đo được & Tài liệu nói\\
\midrule
MCU đạt được bao nhiêu so với ``1 MAC/chu kỳ''? & \expMcuEff\% & PLAN giả định 100\%\\
FPGA nhanh hơn CPU Cortex-A bao nhiêu lần? & \expRatioFpgaPi× (so với Pi 4) & ≈27× (so với A9, Arucu \& Iliev)\\
Độ trễ xấu nhất của Linux PREEMPT\_RT & \expLinuxRtMax{} µs & 30–80 µs (AC6)\\
Độ trễ xấu nhất của Linux thường & \expLinuxStdMax{} µs & vài ms (AC6)\\
Độ trễ ngắt MCU bare-metal & \expMcuIrqRange{} ns & 1–3 µs (AC6)\\
\bottomrule
\end{tabularx}
\par\medskip
\textbf{Quan sát phụ:} một lõi CPU x86 desktop lọc được khoảng 140 MSPS với bài 16-tap.\\
Vậy FPGA không thắng nhờ tốc độ thô ở bài nhỏ, mà nhờ \textbf{thông lượng giữ nguyên khi tăng số tap}, \textbf{độ trễ cố định} và \textbf{công suất thấp}.
\end{frame}

% ---------------------------------------------------------------- 10
\begin{frame}{\MCU: ECU ô tô}
\begin{columns}[T]
\begin{column}{0.55\linewidth}\footnotesize
\begin{itemize}
\item Xe phổ thông có 25–35 MCU, xe sang 60–70
\item AURIX TC3xx: 4/6 lõi chạy \textbf{lockstep}, nhắm tới \textbf{ASIL-D} (ISO 26262), dùng cho powertrain và inverter xe điện
\item AEC-Q100 Grade 0: chạy được từ \textbf{−40 đến 150~°C}; NXP S32K14xW cam kết cung ứng 15 năm
\item Khởi động vài ms, ngắt dưới 1 µs, giá vài USD
\end{itemize}
\end{column}
\begin{column}{0.42\linewidth}\small
\soquyetdinh{\textbf{Vì sao không chọn \SoC:} Linux khởi động mất vài giây, độ trễ không xác định}
\soquyetdinh{\textbf{Vì sao không chọn \FPGA:} đắt; cấu hình nằm trong SRAM nên dễ bị lỗi bức xạ (SEU); khó chứng nhận cho hàng triệu xe}
\soquyetdinh{\textbf{Ưu thế dừng ở:} ADAS và giải trí dùng SoC; xu hướng domain controller gom nhiều ECU thành ít bộ hơn}
\end{column}
\end{columns}
\vfill\scriptsize\textcolor{cInk2}{Cũng chiếm ưu thế: điều khiển động cơ FOC (dưới 21 µs/50 µs), cảm biến chạy pin (300 nA, hơn 10 năm trên pin CR2032), đồ gia dụng}
\end{frame}

% ---------------------------------------------------------------- 11
\begin{frame}{\SoC: router Wi-Fi và điện thoại}
\begin{columns}[T]
\begin{column}{0.55\linewidth}\footnotesize
\begin{itemize}
\item \textbf{Router:} Qualcomm IPQ gom trên một chip CPU Cortex-A chạy Linux/OpenWrt, radio Wi-Fi, switch Ethernet, \textbf{khối xử lý gói tin cứng} (37,5 Mpps) và khối mã hóa
\item \textbf{Điện thoại:} Snapdragon 8 Elite gom CPU, GPU, NPU, ISP và modem 5G trên một chip 3 nm
\item CPU chạy Linux lo phần điều khiển (control plane); khối cứng lo phần chuyển dữ liệu (data plane)
\end{itemize}
\end{column}
\begin{column}{0.42\linewidth}\small
\soquyetdinh{\textbf{Vì sao không chọn \MCU:} không có MMU nên không chạy Linux; phải tự viết lại TCP/IP stack}
\soquyetdinh{\textbf{Vì sao không chọn \FPGA:} viết cả network stack và giao diện bằng HDL quá đắt cho thiết bị tiêu dùng}
\soquyetdinh{\textbf{Ưu thế dừng ở:} điều khiển thời gian thực cứng, nên trong máy vẫn có thêm MCU nhỏ}
\end{column}
\end{columns}
\vfill\scriptsize\textcolor{cInk2}{Cũng chiếm ưu thế: camera IP, TV và set-top box, gateway IoT, giải trí trên xe}
\end{frame}

% ---------------------------------------------------------------- 12
\begin{frame}{\FPGA: dòng dữ liệu song song và tạo mẫu chip}
\footnotesize\textcolor{cInk2}{``Tính toán song song'' nói chung thì GPU phổ biến hơn. FPGA thắng khi dữ liệu \textbf{chảy liên tục từ ADC}, cần \textbf{độ trễ cố định}, và chuẩn còn thay đổi hoặc số lượng không đủ lớn để làm ASIC.}
\begin{columns}[T]
\begin{column}{0.55\linewidth}\footnotesize
\begin{itemize}
\item \textbf{Siêu âm y tế:} beamforming 128–1024 kênh, 4,83 GS/s trên một FPGA; nối thẳng với front-end analog
\item \textbf{Trạm 5G (radio unit):} low-PHY và DFE trên RFSoC hoặc Agilex; chuẩn O-RAN còn đổi nên cần nạp lại được
\item \textbf{Tạo mẫu ASIC:} HAPS ghép nhiều FPGA chạy RTL tới 1,6 tỉ cổng trước khi đúc chip
\item \textbf{HFT:} tick-to-trade 13,9 ns
\end{itemize}
\end{column}
\begin{column}{0.42\linewidth}\small
\soquyetdinh{\textbf{Vì sao không chọn \MCU:} thiếu ≈3,5× năng lực dù đạt hiệu suất 100\%}
\soquyetdinh{\textbf{Vì sao không chọn \SoC:} có thể đủ thông lượng nhờ SIMD/GPU, nhưng mất tính xác định và tốn điện}
\soquyetdinh{\textbf{Ưu thế dừng ở:} số lượng lớn và thuật toán ổn định thì ASIC thắng (ISP, modem trong điện thoại)}
\end{column}
\end{columns}
\end{frame}

% ---------------------------------------------------------------- 13
\begin{frame}{Phương án lai CPU + FPGA: ADAS và điều khiển công nghiệp}
\begin{itemize}
\item Bài toán cần \textbf{cả} hệ điều hành phong phú \textbf{lẫn} timing hoặc xử lý song song cứng
\item Zynq UltraScale+ gom trên một chip: 4×A53 chạy Linux, 2×R5 lockstep chạy real-time, và fabric FPGA
\item Lecture 3, Option \#2: CPU lo tính toán mềm dẻo, FPGA giữ timing actuator kể cả khi CPU bận
\item FIR trên SoC FPGA+ARM nhanh gấp 8,86× ARM đơn, giá/hiệu năng tốt nhất trong các cấu hình được thử \nhan{nguồn}
\end{itemize}
\vfill
\soquyetdinh{\textbf{Cái giá:} chi phí cao hơn và linh hoạt thấp hơn so với một CPU đơn}
\end{frame}

% ---------------------------------------------------------------- 14
\begin{frame}{Tóm tắt các lĩnh vực đặc trưng}
\small
\begin{tabularx}{\linewidth}{@{}l Y Y Y@{}}
\toprule
& Lĩnh vực đặc trưng & Yếu tố quyết định & Vì sao không chọn hai nền tảng kia\\
\midrule
\MCU  & ECU ô tô, động cơ, cảm biến chạy pin & ASIL-D, 150~°C, dòng nghỉ cỡ nA, vài USD & SoC chậm khởi động và không xác định; FPGA đắt\\
\SoC  & Router Wi-Fi, điện thoại, camera IP & Linux + khối tăng tốc cứng & MCU không có MMU; FPGA tốn công HDL\\
\FPGA & Siêu âm, 5G radio, tạo mẫu ASIC, HFT & Song song + I/O trực tiếp + tái cấu hình & MCU thiếu sức; SoC không xác định\\
Lai   & ADAS, công nghiệp & OS + timing cứng & Một nền tảng đơn không đủ\\
\bottomrule
\end{tabularx}
\end{frame}

% ---------------------------------------------------------------- 15
\begin{frame}{Ma trận quyết định}
\small
\begin{tabularx}{\linewidth}{@{}Y l Y@{}}
\toprule
Nếu bài toán\ldots & Chọn & Vì\\
\midrule
Chạy pin nhiều năm, tác vụ nhẹ & \MCU & 300 nA ở Standby\\
Deadline vài chục µs, cần chứng nhận an toàn & \MCU & Ngắt dưới 1 µs; ASIL-D\\
Cần Linux, camera, mạng; timing mềm & \SoC & CoreMark gấp ≈16×; driver có sẵn\\
Trên ≈0,5 G MAC/s liên tục, I/O tốc độ cao & \FPGA{} / lai & 1 mẫu/chu kỳ, ≈48 G MAC/s\\
Độ trễ cỡ ns & \FPGA & 13,9 ns\\
Cần cả OS lẫn timing cứng & Lai & Zynq, Lecture 3\\
\bottomrule
\end{tabularx}
\par\medskip\centering
Muốn timing và thông lượng cao hơn thì phải trả bằng \textbf{giá cao hơn} và \textbf{linh hoạt thấp hơn}
\end{frame}

% ---------------------------------------------------------------- 16
\begin{frame}{Hạn chế}
\begin{itemize}
\item Số liệu tài liệu đến từ nhiều nguồn với điều kiện đo khác nhau
\item Thí nghiệm của nhóm chỉ có \textbf{một} bài toán (FIR 16-tap); ba board ở ba phân khúc giá khác nhau
\item Công suất FPGA là \textbf{ước lượng của Vivado}, không phải số đo
\item Độ trễ ở ba nền tảng là ba đại lượng khác nhau (ngắt, lập lịch, pipeline)
\item Số liệu thị phần đến từ các công ty nghiên cứu thị trường và có thể lệch nhau
\end{itemize}
\end{frame}

% ---------------------------------------------------------------- 17
\begin{frame}[plain]
\vfill\centering
{\Huge\bfseries Không có nền tảng tốt nhất}\par\bigskip
{\Large chỉ có nền tảng \textbf{đủ} cho con số quyết định của bài toán,\\ với chi phí thấp nhất}\par\vfill
{\color{cInk2}Hỏi đáp}
\vfill
\end{frame}

% ---------------------------------------------------------------- phu luc
\appendix
\begin{frame}{Phụ lục: chuẩn bị cho câu hỏi}
\small
\begin{itemize}
\item \textbf{``MCU đã tối ưu chưa?''} Có bản CMSIS-DSP. M3 không có SIMD, bản tốt nhất đạt \expMcuBestCps{} chu kỳ/mẫu.
\item \textbf{``Jitter của FPGA bằng 0 thật à?''} Mô phỏng RTL không có thời gian thật. Trên chip, độ trễ cố định theo số chu kỳ, chỉ còn clock jitter cỡ ps.
\item \textbf{``Sao so chip với board?''} Chỉ để so bậc độ lớn, và nhóm đã ghi rõ trên từng bảng.
\item \textbf{``Công suất FPGA đo bằng gì?''} Không đo; đó là Vivado \texttt{report\_power} (vector-less) \nhan{ước lượng công cụ}.
\item \textbf{``Sao không có NRE?''} Chưa có số công bố. Nhóm dùng thời gian thực tế viết bản C và bản HDL làm dẫn chứng định tính.
\item \textbf{``Tăng lên 64 tap thì sao?''} MCU chậm đi khoảng 4×. FPGA vẫn 1 mẫu/chu kỳ, chỉ tốn thêm DSP và vài tầng pipeline.
\end{itemize}
\end{frame}

\begin{frame}{Nguồn chính}
\scriptsize
\begin{itemize}
\item ST: STM32H7/F4/MP1 overview, STM32H743; OpenBenchmarking: Raspberry Pi 4
\item Arucu \& Iliev, J. Low Power Electron. Appl., 2025 (FIR trên FPGA, GPU, CPU)
\item AC6: độ trễ Linux PREEMPT\_RT; A-Team Insight: Exegy/AMD STAC-T0
\item Embedded World 02/2026 (số MCU mỗi xe); Automotive World (AURIX TC3xx); Future Electronics (NXP S32K Grade 0)
\item Macnica, CNX Software (Qualcomm IPQ); GSMArena (Snapdragon 8 Elite)
\item Kou et al., 2023 (beamforming siêu âm trên FPGA); HPCwire (Fujitsu + Xilinx O-RAN); Intel (Agilex O-RU); Synopsys HAPS
\item eeNews Europe (Zynq UltraScale+); Lecture 2, Lecture 3; Lee \& Seshia mục 8.1
\end{itemize}
\par\medskip
Danh sách đầy đủ kèm đường link nằm trong báo cáo (\texttt{report.pdf}).
\end{frame}

\end{document}
